% Part: normal-modal-logic
% Chapter: tableaux
% Section: soundness

\documentclass[../../../include/open-logic-section]{subfiles}

\begin{document}

\olfileid{nml}{tab}{sou}

\olsection{درستی برای~\Log{K}}

\begin{editorial}
  این اثبات درستی، اثبات درستیِ منطق گزاره‌ای کلاسیک را بازبه‌کار
  می‌گیرد؛ یعنی همه‌چیز را از آغاز اثبات می‌کند. برای اثباتی
  خودبسنده مناسب است. اگر درستیِ تابلوی معمول را پیش‌تر دیده باشید،
  تکراری خواهد بود. در نظر است امکان انتخاب میان نسخهٔ خودبسنده و
  نسخه‌ای مبتنی بر حالت غیرموجهاتی فراهم شود.
\end{editorial}

\begin{explain}
  برای نشان‌دادن درستیِ !!{tableau}های پیشونددار باید نشان دهیم اگر
  \[
  \sFmla{\True}{!B_1}[1], \dots, \sFmla{\True}{!B_n}[1], \sFmla{\False}{!A}[1]
  \]
  یک !!{tableau} بسته داشته باشد، آنگاه~$!B_1, \dots, !B_n \Entails
  !A$. اثبات عکس نقیض آسان‌تر است: اگر برای مدل~$\mModel{M}$ و
  جهان~$w$ای، $\mSat{M}{!B_i}[w]$ برای همهٔ~$i=1$، \dots،~$n$، اما
  $\mSat/{M}{!A}[w]$، آنگاه هیچ !!{tableau}ای نمی‌تواند بسته شود. چنین
  پادمدلی نشان می‌دهد فرض‌های آغازین !!{tableau} ارضاپذیرند. راهبرد
  اثبات این است که نشان دهیم هرگاه همهٔ !!{formula}های پیشونددار روی
  شاخه‌ای از !!a{tableau} ارضاپذیر باشند، هر اعمال قاعده دست‌کم یک
  شاخهٔ گسترش‌یافته به دست می‌دهد که آن نیز ارضاپذیر است. چون شاخه‌های
  بسته ارضاناپذیرند، هر !!{tableau} برای مجموعه‌ای ارضاپذیر از
  !!{formula}های پیشونددار باید دست‌کم یک شاخهٔ باز داشته باشد.

  برای به‌کارگیری این راهبرد در حالت موجهاتی، باید تعریف «ارضاپذیر» را
  به مدل‌های موجهاتی و پیشوندها گسترش دهیم. پس از این کار، اثبات
  سرراست است.
\end{explain}

\begin{defn}
  بگذارید~$P$ مجموعه‌ای از پیشوندها باشد؛ یعنی~$P \subseteq
  (\PosInt)^* \setminus \{\emptyseq\}$، و~$\mModel{M}$ مدلی باشد.
  تابع~$f\colon P \to W$ یک \emph{تفسیرِ}~$P$ در~$\mModel{M}$ است
  اگر هرگاه~$\sigma$ و~$\sigma.n$ هر دو در~$P$ باشند، آنگاه
  $Rf(\sigma)f(\sigma.n)$.

  نسبت به یک تفسیر از پیشوندهای~$P$ می‌توانیم چنین تعریف کنیم:
  \begin{enumerate}
  \item $\mModel{M}$، $\sFmla{\True}{!A}[\sigma]$ را ارضا می‌کند اگر
    و تنها اگر~$\mSat{M}{!A}[f(\sigma)]$.
  \item $\mModel{M}$، $\sFmla{\False}{!A}[\sigma]$ را ارضا می‌کند
    اگر و تنها اگر~$\mSat/{M}{!A}[f(\sigma)]$.
  \end{enumerate}
\end{defn}

\begin{defn}
  بگذارید~$\Gamma$ مجموعه‌ای از !!{formula}های پیشونددار باشد و
  $P(\Gamma)$ مجموعهٔ پیشوندهای ظاهرشونده در آن. اگر~$f$ تفسیری از
  $P(\Gamma)$ در~$\mModel{M}$ باشد، می‌گوییم~$\mModel{M}$، $\Gamma$
  را نسبت به~$f$ ارضا می‌کند، $\mSat{M}{\Gamma}[f]$، اگر~$\mModel{M}$
  هر !!{formula} پیشونددار در~$\Gamma$ را نسبت به~$f$ ارضا کند.
  $\Gamma$ \emph{ارضاپذیر} است اگر و تنها اگر مدل~$\mModel{M}$ و
  تفسیر~$f$ از~$P(\Gamma)$ وجود داشته باشند چنان‌که
  $\mSat{M}{\Gamma}[f]$.
\end{defn}

\begin{prop}
  اگر~$\Gamma$، هم~$\sFmla{\True}{!A}[\sigma]$ و هم
  $\sFmla{\False}{!A}[\sigma]$ را برای یک !!{formula}~$!A$ و
  پیشوند~$\sigma$ شامل شود، آنگاه~$\Gamma$ ارضاناپذیر است.
\end{prop}

\begin{proof}
  هیچ مدل~$\mModel{M}$ و تفسیر~$f$ از~$P(\Gamma)$ نمی‌تواند وجود
  داشته باشد که هم~$\mSat{M}{!A}[f(\sigma)]$ و هم
  $\mSat/{M}{!A}[f(\sigma)]$.
\end{proof}

\begin{thm}[درستی]
  \ollabel{thm:tableau-soundness}
  اگر~$\Gamma$ یک !!{tableau} بسته داشته باشد، $\Gamma$ ارضاناپذیر است.
\end{thm}

\begin{proof}
شاخه‌ای از !!a{tableau} را ارضاپذیر می‌نامیم اگر و تنها اگر مجموعهٔ
!!{signed formula}های روی آن ارضاپذیر باشد؛ و !!a{tableau} را ارضاپذیر
می‌نامیم اگر دست‌کم یک شاخهٔ ارضاپذیر داشته باشد.

نشان می‌دهیم: گسترش یک !!{tableau} ارضاپذیر با یکی از قواعد استنتاج
همواره !!{tableau}ای ارضاپذیر به دست می‌دهد. این قضیه را ثابت می‌کند:
هر !!{tableau} بسته با اعمال قواعد استنتاج بر !!{tableau}ای حاصل
می‌شود که تنها از فرض‌های~$\Gamma$ تشکیل شده است. پس اگر~$\Gamma$
ارضاپذیر می‌بود، هر !!{tableau} برای آن ارضاپذیر می‌بود. اما یک
!!{tableau} بسته آشکارا ارضاپذیر نیست، زیرا همهٔ شاخه‌هایش بسته‌اند و
شاخه‌های بسته ارضاناپذیرند.

فرض کنید یک !!{tableau} ارضاپذیر، یعنی !!a{tableau}ای با دست‌کم یک
شاخهٔ ارضاپذیر، داریم. اعمال یک قاعدهٔ استنتاج یا !!{signed formula}هایی
به شاخه می‌افزاید یا شاخه را به دو شاخه می‌شکافد. اگر !!{tableau}
شاخه‌ای ارضاپذیر داشته باشد که اعمال قاعدهٔ مورد نظر آن را گسترش ندهد،
آن شاخه در !!{tableau} گسترش‌یافته ارضاپذیر می‌ماند؛ پس تابلوی
گسترش‌یافته ارضاپذیر است. بنابراین فقط حالتی را بررسی می‌کنیم که قاعده
بر شاخه‌ای ارضاپذیر اعمال شود.

بگذارید~$\Gamma$ مجموعهٔ !!{signed formula}های روی آن شاخه و
$\sFmla{S}{!A}[\sigma] \in \Gamma$ همان !!{signed formula}ای باشد که
قاعده بر آن اعمال می‌شود. اگر قاعده شاخه را نشکند، باید نشان دهیم شاخهٔ
گسترش‌یافته، یعنی~$\Gamma$ همراه با نتایج قاعده، همچنان ارضاپذیر است.
اگر قاعده شاخه را بشکند، باید نشان دهیم دست‌کم یکی از دو شاخهٔ حاصل
ارضاپذیر است.
نخست استنتاج‌های ممکن با تنها یک مقدمه را بررسی می‌کنیم.
\begin{enumerate}
\item شاخه با اعمال~$\TRule{\True}{\lnot}$ بر
  $\sFmla{\True}{\lnot !B}[\sigma] \in \Gamma$ گسترش می‌یابد. آنگاه
  شاخهٔ گسترش‌یافته !!{signed formula}های~$\Gamma \cup
  \{\sFmla{\False}{!B}[\sigma]\}$ را شامل می‌شود. فرض کنید
  $\mSat{M}{\Gamma}[f]$. به‌ویژه~$\mSat{M}{\lnot !B}[f(\sigma)]$؛ پس
  $\mSat/{M}{!B}[f(\sigma)]$، یعنی~$\mModel{M}$ نسبت به~$f$،
  $\sFmla{\False}{!B}[\sigma]$ را ارضا می‌کند.
\item شاخه با اعمال~$\TRule{\False}{\lnot}$ بر
  $\sFmla{\False}{\lnot !B}[\sigma] \in \Gamma$ گسترش می‌یابد: تمرین.
\item شاخه با اعمال~$\TRule{\True}{\land}$ بر
  $\sFmla{\True}{!B \land !C}[\sigma] \in \Gamma$ گسترش می‌یابد، که
  دو !!{signed formula} تازه روی شاخه به دست می‌دهد:
  $\sFmla{\True}{!B}[\sigma]$ و~$\sFmla{\True}{!C}[\sigma]$. فرض کنید
  $\mSat{M}{\Gamma}[f]$، و به‌ویژه
  $\mSat{M}{!B \land !C}[f(\sigma)]$. آنگاه
  $\mSat{M}{!B}[f(\sigma)]$ و~$\mSat{M}{!C}[f(\sigma)]$. یعنی
  $\mModel{M}$ نسبت به~$f$، هم~$\sFmla{\True}{!B}[\sigma]$ و هم
  $\sFmla{\True}{!C}[\sigma]$ را ارضا می‌کند.
\item شاخه با اعمال~$\TRule{\False}{\lor}$ بر
  $\sFmla{\False}{!B \lor !C}[\sigma] \in \Gamma$ گسترش می‌یابد: تمرین.
\item شاخه با اعمال~$\TRule{\False}{\lif}$ بر
  $\sFmla{\False}{!B \lif !C}[\sigma] \in \Gamma$ گسترش می‌یابد. دو
  !!{signed formula} تازه روی شاخه به دست می‌آید:
  $\sFmla{\True}{!B}[\sigma]$ و~$\sFmla{\False}{!C}[\sigma]$. فرض کنید
  $\mSat{M}{\Gamma}[f]$، و به‌ویژه
  $\mSat/{M}{!B \lif !C}[f(\sigma)]$. آنگاه
  $\mSat{M}{!B}[f(\sigma)]$ و~$\mSat/{M}{!C}[f(\sigma)]$. یعنی
  $\mModel{M}, f$ هر دو عبارت~$\sFmla{\True}{!B}[\sigma]$ و
  $\sFmla{\False}{!C}[\sigma]$ را ارضا می‌کند.

\iftag{prvBox}{%
\item شاخه با اعمال~$\TRule{\True}{\Box}$ بر
  $\sFmla{\True}{\Box !B}[\sigma] \in \Gamma$ گسترش می‌یابد:
  \iftag{probBox}{تمرین.}{این کار !!{signed formula} تازهٔ
    $\sFmla{\True}{!B}[\sigma.n]$ را برای یک~$\sigma.n \in
    P(\Gamma)$ روی شاخه به دست می‌دهد (زیرا~$\sigma.n$ باید
    استفاده‌شده باشد). فرض کنید~$\mSat{M}{\Gamma}[f]$، و به‌ویژه
    $\mSat{M}{\Box !B}[f(\sigma)]$. چون~$f$ تفسیری از پیشوندهاست و
    هر دو~$\sigma$ و~$\sigma.n \in P(\Gamma)$، می‌دانیم
    $Rf(\sigma)f(\sigma.n)$. پس~$\mSat{M}{!B}[f(\sigma.n)]$؛ یعنی
    $\mModel{M}, f$، $\sFmla{\True}{!B}[\sigma.n]$ را ارضا می‌کند.}

\item شاخه با اعمال~$\TRule{\False}{\Box}$ بر
  $\sFmla{\False}{\Box !B}[\sigma] \in \Gamma$ گسترش می‌یابد:
  \iftag{probBox}{تمرین.}{این کار !!{signed formula} تازهٔ
    $\sFmla{\False}{!B}[\sigma.n]$ را به دست می‌دهد، که~$\sigma.n$
    پیشوندی تازه روی شاخه است؛ یعنی~$\sigma.n \notin P(\Gamma)$.
    چون~$\Gamma$ ارضاپذیر است، مدل~$\mModel{M}$ و تفسیر~$f$ از
    $P(\Gamma)$ وجود دارند چنان‌که~$\mSat{M}{\Gamma}[f]$، و به‌ویژه
    $\mSat/{M}{\Box !B}[f(\sigma)]$. باید نشان دهیم
    $\Gamma \cup \{\sFmla{\False}{!B}[\sigma.n]\}$ ارضاپذیر است.
    برای این کار، تفسیری از~$P(\Gamma) \cup \{\sigma.n\}$ چنین تعریف
    می‌کنیم:

  چون~$\mSat/{M}{\Box !B}[f(\sigma)]$، جهان~$w \in W$ای وجود دارد که
  $Rf(\sigma)w$ و~$\mSat/{M}{!B}[w]$. بگذارید~$f'$ مانند~$f$ باشد،
  جز آنکه~$f'(\sigma.n) = w$. چون~$f'(\sigma) = f(\sigma)$ و
  $Rf(\sigma)w$، داریم~$Rf'(\sigma)f'(\sigma.n)$؛ پس~$f'$ تفسیری از
  $P(\Gamma) \cup \{\sigma.n\}$ است. آشکارا
  $\mSat/{M}{!B}[f'(\sigma.n)]$. چون برای همهٔ پیشوندهای
  $\sigma' \in P(\Gamma)$، $f(\sigma') = f'(\sigma')$، داریم
  $\mSat{M}{\Gamma}[f']$. پس~$\mModel{M}, f'$،
  $\Gamma \cup \{\sFmla{\False}{!B}[\sigma.n]\}$ را ارضا می‌کند.}
  }{}

\iftag{prvDiamond}{%
\item شاخه با اعمال~$\TRule{\True}{\Diamond}$ بر
  $\sFmla{\True}{\Diamond !B}[\sigma] \in \Gamma$ گسترش می‌یابد:
  \iftag{probDiamond}{تمرین.}{این کار !!{signed formula} تازهٔ
    $\sFmla{\True}{!B}[\sigma.n]$ را به دست می‌دهد، که~$\sigma.n$
    پیشوندی تازه روی شاخه است؛ یعنی~$\sigma.n \notin P(\Gamma)$.
    چون~$\Gamma$ ارضاپذیر است، مدل~$\mModel{M}$ و تفسیر~$f$ از
    $P(\Gamma)$ وجود دارند چنان‌که~$\mSat{M}{\Gamma}[f]$، و به‌ویژه
    $\mSat{M}{\Diamond !B}[f(\sigma)]$. باید نشان دهیم
    $\Gamma \cup \{\sFmla{\True}{!B}[\sigma.n]\}$ ارضاپذیر است.
    برای این کار، تفسیری از~$P(\Gamma) \cup \{\sigma.n\}$ چنین تعریف
    می‌کنیم:

  چون~$\mSat{M}{\Diamond !B}[f(\sigma)]$، جهان~$w \in W$ای وجود دارد
  که~$Rf(\sigma)w$ و~$\mSat{M}{!B}[w]$. بگذارید~$f'$ مانند~$f$ باشد،
  جز آنکه~$f'(\sigma.n) = w$. چون~$f'(\sigma) = f(\sigma)$ و
  $Rf(\sigma)w$، داریم~$Rf'(\sigma)f'(\sigma.n)$؛ پس~$f'$ تفسیری از
  $P(\Gamma) \cup \{\sigma.n\}$ است. آشکارا
  $\mSat{M}{!B}[f'(\sigma.n)]$. چون برای همهٔ پیشوندهای
  $\sigma' \in P(\Gamma)$، $f(\sigma') = f'(\sigma')$، داریم
  $\mSat{M}{\Gamma}[f']$. پس~$\mModel{M}, f'$،
  $\Gamma \cup \{\sFmla{\True}{!B}[\sigma.n]\}$ را ارضا می‌کند.}
\item شاخه با اعمال~$\TRule{\False}{\Diamond}$ بر
  $\sFmla{\False}{\Diamond !B}[\sigma] \in \Gamma$ گسترش می‌یابد:
  \iftag{probDiamond}{تمرین.}{این کار !!{signed formula} تازهٔ
    $\sFmla{\False}{!B}[\sigma.n]$ را برای یک~$\sigma.n \in
    P(\Gamma)$ روی شاخه به دست می‌دهد (زیرا~$\sigma.n$ باید
    استفاده‌شده باشد). فرض کنید~$\mSat{M}{\Gamma}[f]$، و به‌ویژه
    $\mSat/{M}{\Diamond !B}[f(\sigma)]$. چون~$f$ تفسیری از پیشوندهاست
    و هر دو~$\sigma$ و~$\sigma.n \in P(\Gamma)$، می‌دانیم
    $Rf(\sigma)f(\sigma.n)$. پس~$\mSat/{M}{!B}[f(\sigma.n)]$؛ یعنی
    $\mModel{M}, f$، $\sFmla{\False}{!B}[\sigma.n]$ را ارضا می‌کند.}
  }{}
\end{enumerate}
اکنون استنتاج‌های ممکن با دو مقدمه را بررسی می‌کنیم.
\begin{enumerate}
\item شاخه با اعمال~$\TRule{\False}{\land}$ بر
  $\sFmla{\False}{!B \land !C}[\sigma] \in \Gamma$ گسترش می‌یابد و
  دو شاخه به دست می‌دهد: شاخهٔ چپ از
  $\sFmla{\False}{!B}[\sigma]$ ادامه می‌یابد و شاخهٔ راست از
  $\sFmla{\False}{!C}[\sigma]$. فرض کنید~$\mSat{M}{\Gamma}[f]$، و
  به‌ویژه~$\mSat/{M}{!B \land !C}[f(\sigma)]$. آنگاه یا
  $\mSat/{M}{!B}[f(\sigma)]$ یا~$\mSat/{M}{!C}[f(\sigma)]$. در حالت
  نخست~$\mModel{M}, f$، $\sFmla{\False}{!B}[\sigma]$ را ارضا می‌کند؛
  یعنی شاخهٔ چپ ارضاپذیر است. در حالت دوم~$\mModel{M}, f$،
  $\sFmla{\False}{!C}[\sigma]$ را ارضا می‌کند؛ یعنی شاخهٔ راست
  ارضاپذیر است.
\item شاخه با اعمال~$\TRule{\True}{\lor}$ بر
    $\sFmla{\True}{!B \lor !C}[\sigma] \in \Gamma$ گسترش می‌یابد: تمرین.
\item شاخه با اعمال~$\TRule{\True}{\lif}$ بر
    $\sFmla{\True}{!B \lif !C}[\sigma] \in \Gamma$ گسترش می‌یابد: تمرین.
\end{enumerate}
\end{proof}

\begin{prob}
اثبات \olref[nml][tab][sou]{thm:tableau-soundness} را کامل کنید.
\end{prob}

\begin{cor}
\ollabel{cor:entailment-soundness}
اگر~$\Gamma \Proves !A$، آنگاه~$\Gamma \Entails !A$.
\end{cor}

\begin{proof}
  اگر~$\Gamma \Proves !A$، آنگاه برای یک~$!B_1$، \dots، $!B_n \in
  \Gamma$، مجموعهٔ~$\Delta = \{\sFmla{\False}{!A}[1],
  \sFmla{\True}{!B_1}[1], \dots, \sFmla{\True}{!B_n}[1]\}$ یک
  !!{tableau} بسته دارد. می‌خواهیم نشان دهیم~$\Gamma \Entails !A$.
  خلاف آن را فرض کنید؛ پس برای یک~$\mModel{M}$ و~$w$،
  $\mSat{M}{!B_i}[w]$ برای~$i=1$، \dots،~$n$، اما
  $\mSat/{M}{!A}[w]$. بگذارید~$f(1) = w$؛ آنگاه~$f$ تفسیری از
  $P(\Delta)$ در~$\mModel{M}$ است و~$\mModel{M}$، $\Delta$ را نسبت
  به~$f$ ارضا می‌کند. اما بنا بر \olref{thm:tableau-soundness}،
  $\Delta$ ارضاناپذیر است، زیرا یک !!{tableau} بسته دارد؛ تناقض. پس
  سرانجام باید~$\Gamma \Entails !A$.
\end{proof}

\begin{cor}
\ollabel{cor:weak-soundness}
اگر~$\Proves !A$، آنگاه~$!A$ در همهٔ مدل‌ها صادق است.
\end{cor}

\end{document}
